% Phase 11 task 11.1: abstract skeleton. Full draft lands with the
% writing pass that follows the claim-to-experiment mapping in 11.2.
\todo{draft abstract (200--250 words).}

Multi-agent LLM systems that share a vector memory are vulnerable to
memory-poisoning attacks: a single compromised agent can inject
retrieval-time triggers that steer every downstream agent. We introduce
\sys{}, a runtime monitoring layer that combines (i) six per-agent
behavioral signals, (ii) a Bayesian Online Change-Point Detector (BOCPD)
on the dominant signal, (iii) a Chronos-2-Small probabilistic forecaster
that supplies a temporal prior on future agent behavior, and (iv) a
health-weighted decentralized task allocator that routes work away from
drifting agents. We evaluate \sys{} on two memory-poisoning attacks
(MINJA, AgentPoison) across six experimental configurations (full
system, three ablations, and two reactive baselines). We report Advance
Warning Time (\awt{}), allocation efficiency under attack, and ROC
characteristics, and map each number back to the five claims C1--C5
listed in \cref{sec:experiments}.
